\section{三纤维结构：尺度、相位与素数（亦称三寄存器）}

本文将“三纤维”（亦称“三寄存器”）严格化为三条先验纤维：尺度（外表）纤维、相位纤维与素数纤维（prime-register）。其共同作用是：在保证局域归一化、可逆回放与核验闭环的同时，给出足以表达算术与谱读出的最小可见参数集合。

\subsection{外表纤维与尺度纤维}
\begin{definition}[外表纤维 $Z$]\label{def:Z-fiber}
外表纤维取 Zeckendorf 稳定地址空间
$$
Z := \Bigl\{(c_k)_{k\ge 1}\in\{0,1\}^{\NN}:\ \forall k,\ c_kc_{k+1}=0,\ \text{且仅有限个 }1\Bigr\}.
$$
令 Fibonacci 数列满足 $F_1=F_2=1$，$F_{k+2}=F_{k+1}+F_k$。定义值映射
$$
\mathrm{Val}(c) := \sum_{k\ge 1} c_k F_{k+1}.
$$
则 $\mathrm{Val}$ 将 $Z$ 与 $\NN$ 同构，并为每个值类选取唯一规范代表。
\end{definition}

\begin{definition}[尺度纤维 $S_\varphi$]\label{def:scale-fiber}
对无理数 $\alpha\in(0,1)$ 的连分数展开 $\alpha=[0;a_1,a_2,\dots]$，Ostrowski 表示 \cite{Ostrowski1922} 给出唯一数字序列 $(d_n)_{n\ge 1}$ 满足合法约束并表示
$$
N=\sum_{n\ge 1} d_n q_{n-1},
$$
其中 $q_n$ 为收敛分母。取 $\alpha=\varphi^{-1}$（$\varphi=(1+\sqrt{5})/2$）时，合法约束退化为 Zeckendorf 语法（禁止相邻 $11$）：有限分辨率空间为 $X_m$，逆极限为 $X_\infty$。尺度纤维定义为该稳定语法及其折叠归一化 $\Fold_m,\Fold_\infty$ 所构成的地址层。
\end{definition}

尺度纤维的意义在于：它把“局域归一化”固定为可终止、可合流、正规形唯一的重写系统，从而为算术涌现与可见寄存器语义提供可计算的规范代表（见附录 \ref{app:fold-rewriting}）。

\subsection{相位纤维}
\begin{definition}[相位纤维 $T$]\label{def:phase-fiber}
设相空间为紧阿贝尔群（例如 $\TT^n$），其角色基为
$$
\chi_k(x)=e^{2\pi i\langle k,x\rangle},\qquad k\in\ZZ^n.
$$
相位纤维 $T$ 指由角色结构诱导的谱坐标层，用于将可加观测量提升为可乘权重并实现卷积对角化；其在后续的 unitary slice 约束下表现为可复核的相位参数轴 $\theta$。
\end{definition}
相位纤维提供“连续可见参数”的最小化：在本文闭包中，可核对的连续参数被锁定为 $u=e^{i\theta}$ 的相位，而不包含径向自由度。

\subsection{素数纤维（prime-register）}
\begin{definition}[素数纤维 $P$（prime-register）]\label{def:prime-register-fiber}
令 $\mathcal{P}$ 为素数集合，定义有限支撑指数向量空间
$$
P := \NN^{(\mathcal{P})}=\Bigl\{r:\mathcal{P}\to\NN\ \Bigm|\ r(p)=0\ \text{除有限多个 }p\Bigr\}.
$$
定义值映射
$$
N(r)=\prod_{p\in\mathcal{P}} p^{r(p)}.
$$
由算术基本定理，$N(\cdot)$ 给出 $P$ 与 $\NN_{\ge 1}$ 的同构，从而把乘法线性化为坐标加法。由于每个状态仅激活有限多素数轴，该纤维满足“无限坐标系 + 局部触发”的可实现约束。
\end{definition}

此外，素数寄存器纤维同时承担“残差承载”角色：折叠的纤维差异可嵌入有限支撑素数幂坐标，使不可逆压缩被改写为跨纤维存储，从而实现无损提升（见附录 \ref{app:prime-uplift}）。
在局域并行设定下，“加法轨道寄存器（轨道残差）”与“乘法素数寄存器（唯一分解）”的作用等价性见附录 \ref{app:bridge-rh-obligations} 的命题 \ref{prop:orbit-register-eq-prime-register}；事件编码 $\mathrm{code}/\mathrm{decode}$ 的完全形式化与 unitary slice 核验的数值证书构造见附录 \ref{app:inject-code-audit}。

\subsection{三寄存器闭包与最小性问题}
三纤维的组合解释了 $(\varphi,e,\pi)$ 等常数在不同层级的出现：$\varphi$ 绑定于尺度纤维的计数斜率与稳定语法；$e$ 与 $2\pi$ 绑定于相位纤维的角色归一化；素数结构绑定于素数寄存器纤维的坐标化与完备化轴。更一般地，三寄存器闭包的最小性可表述为：若要求同时满足
\begin{enumerate}[label=(\roman*), ref=\roman*]
  \item 局域归一化可执行（有限窗口/有限状态）；
  \item 归一化可逆回放（信息不湮灭）；
  \item 连续参数可核验（可见轴可由证书验证），
\end{enumerate}
则可见自由度必须受限于尺度的离散规范代表、相位的紧致参数与素数残差的离散承载；任何额外非紧致连续自由度将表现为额外寄存器需求并破坏类型闭包（见第 \ref{sec:fourth-register} 节）。

\subsection{类型化纤维演算：排序、索引与资源上下文}\label{subsec:typed-fiber-calculus}
本文不把“类型”理解为编程语言语法，而把它作为\textbf{可核验可读出性（partial readout）}的数学约束：在既定闭包与核验协议下，\emph{可构造}与\emph{可运算}本身就是一个可判定的语义层。\par
核心原则是：\textbf{类型不相容}并非数值误差，而是\textbf{偏函数未定义}；实现层可用显式标记承载未定义，但论文语义以偏函数为准。
基础值域与未定义约定见定义 \ref{def:valdom-undefined}。

\begin{definition}[索引化排序与资源上下文]\label{def:indexed-sorts-resctx}
为将“分辨率/精度/预算”提升为类型索引，引入三类索引化排序：\par
\noindent\textbf{(1) 尺度排序 $\mathsf{Scale}[m]$：}其中 $m\in\NN$ 为离散窗口；其居住项是分辨率 $m$ 下的稳定规范代表（例如 $X_m$、或其等价的折叠正规形）。\par
\noindent\textbf{(2) 相位排序 $\mathsf{Phase}[p]$：}其中 $p\in\NN$ 为可重复核验的精度索引；本文默认以 dyadic 网格 $D_p=\{k/2^p\}$（或等价的有理区间端点体系）作为可复算载体，使 unitary slice 的连续输入可被强类型化为 $(L,\theta,p)\in\Addr_{\mathrm{US}}$（定义 \ref{def:us-typed-address-null}）。\par
\noindent\textbf{(3) 残差排序 $\mathsf{Resid}[\beta]$：}其中 $\beta$ 表示记录轴预算（例如 prime-register 的模积/位数预算，或端点精度载荷预算）；其居住项承载单射化所需的侧信息（见第 \ref{sec:ledger} 节与附录 \ref{app:prime-uplift}）。\par
以资源上下文 $\Gamma:=\mathsf{ResCtx}(m,p,\beta,\dots)$ 记录当前可用的索引与预算，并用类型判定式 $\Gamma\vdash t:\tau$ 表示“在给定资源下 $t$ 为可核验可构造的 $\tau$ 型项”。
\end{definition}

\begin{definition}[共同细化、显式提升与类型失败]\label{def:common-refinement-type-failure}
任何二元运算只在输入被\emph{显式提升}到共同索引且预算覆盖记录成本时定义。更具体地，对同一排序族 $\tau[-]$：若
$$
\Gamma\vdash x:\tau[i],\qquad \Gamma\vdash y:\tau[j],
$$
则 $x\circ y$ 仅在存在可核验提升证书使二者落入共同索引 $k\succeq i,j$（例如 $k=\max(i,j)$）时才被允许；否则该项在语义上未定义并归入 $\NULL$ 机制。\par
因此，“精度不同则不能直接运算”在本文语义中被解释为：\textbf{索引不一致导致项不可构造（偏函数未定义）}，除非提供提升算子（例如外向舍入的区间提升）及其证书化实现。
\end{definition}

\begin{remark}[unitary slice 的类型封闭与第四纤维]\label{rem:unitary-slice-type-closure}
unitary slice 核验协议把连续可见参数轴强类型化为相位 $u=e^{i\theta}$（定义 \ref{def:unitary-slice-audit-axis}）。当论证不可避免地要求一般分解 $u=re^{i\theta}$ 且 $r\neq 1$ 时，类型必须外延为
$$
\mathsf{Phase}\leadsto \mathsf{Phase}\times \mathsf{Radial},
$$
其中 $\mathsf{Radial}$ 作为额外连续自由度的承载即对应“第四纤维（寄存器）”；其结构与定量障碍在第 \ref{sec:fourth-register} 节给出。
\end{remark}

